% ---------------------------------------------------------------------------
% Team 38 Final Report -- Section 4c (Problems encountered) and Section 5
% (Plans for Spring). Author: Dastan Issenov.
%
% Needs in the preamble:
%   \usepackage{longtable}
%   \usepackage{booktabs}
%   \usepackage{array}
% ---------------------------------------------------------------------------

\subsection{Problems Encountered}
\label{sec:problems}

The three main problems were the lack of a GPU, a weak predictive signal and
gaps in the data sources. Most of them were solved or worked around during the
semester; the rest are scheduled for spring. Table~\ref{tab:problems} lists each
problem, what happened and how it was solved or will be solved.

{\small
\begin{longtable}{@{}>{\raggedright\arraybackslash}p{0.17\linewidth}
                    >{\raggedright\arraybackslash}p{0.38\linewidth}
                    >{\raggedright\arraybackslash}p{0.38\linewidth}@{}}
\caption{Problems encountered during the Fall 2026 semester.}
\label{tab:problems}\\
\toprule
\textbf{Problem} & \textbf{What happened} & \textbf{How we solved it, or will solve it} \\
\midrule
\endfirsthead
\multicolumn{3}{@{}l}{\small\emph{Table~\ref{tab:problems} continued}}\\
\toprule
\textbf{Problem} & \textbf{What happened} & \textbf{How we solved it, or will solve it} \\
\midrule
\endhead
\midrule
\multicolumn{3}{r@{}}{\small\emph{continued on the next page}}\\
\endfoot
\bottomrule
\endlastfoot

Late start of implementation &
The first part of the semester went into research and design: data sources,
architecture and the GPU plan. Coding started on 28 October. &
We worked in short phases, with tests and linters before every commit. The
website, the data pipeline and the first trained model were finished within the
semester. \\
\addlinespace

No GPU &
The RTX 5090 has not arrived and NU HPC access has not been requested yet, so
all training ran on an 8-core laptop CPU. &
Models sit behind an interface that swaps placeholder models for real ones by
configuration. Gradient-boosted models were trained on CPU. LLM news sentiment
and the LLM trade reviewer moved to spring; a rule-based reviewer runs
meanwhile. A day-by-day GPU plan is ready. \\
\addlinespace

Scope differs from the requirements document &
The requirements document specified the StockNet benchmark (2014--2015 tweets),
next-day direction and a stacking meta-learner with news sentiment. The built
system uses live data for US and Kazakhstan stocks, a 5-day horizon and an
averaged LightGBM + XGBoost ensemble without sentiment. &
In spring we add news sentiment and run the price-only vs.\ sentiment ablation
from the original plan, including a StockNet benchmark run for comparison with
published results. \\
\addlinespace

Accuracy below a trivial baseline &
The first model reached 52.5\% accuracy, below the 54.6\% of always predicting
``up'', because it mostly guessed the direction of the whole market. &
We changed the label to ``beats the day's median return'', added
cross-sectional rank features and averaged LightGBM with XGBoost. The ranking
signal rose from IC~0.013 to 0.023 ($t = 4.2$), and every metric is now
reported next to a simple baseline. \\
\addlinespace

No profit after costs &
The v2 strategy's Sharpe ratio (1.38) only ties the equal-weight benchmark
(1.36), with higher risk, and a plain 21-day momentum rule does better (1.55). &
The model stays in demo mode under a promotion rule fixed before testing, and
the website shows a ``Demo predictions'' badge. Spring: a beta-neutral label and
volatility-scaled positions. \\
\addlinespace

Survivorship bias &
The stock universe is today's 100 most liquid companies, the winners of recent
years. This inflates backtest returns and flatters momentum rules. &
Documented in the data report. Spring: build the universe from historical index
membership on each date. \\
\addlinespace

Risk of look-ahead bias &
Company filings are often published after the market close and macro data are
revised later, so naive date joins leak future information. &
Each value enters only after its publication time (SEC acceptance timestamp,
first-release macro values). An automated test checks that adding future data
never changes past features. Design choices used 2016--2018 only; 2019--2026
was scored once. \\
\addlinespace

Gaps and quirks in data sources &
FRED serves the high-yield credit spread only from October 2023. Tradernet
blocked requests whose User-Agent held a URL, and returns bars in an unusual
column order with Moscow-midnight timestamps. HTTP logs exposed the FRED key,
and an NBK outage stopped a backfill. &
We removed the credit spread, changed the User-Agent, fixed the column order
and time zone, stopped logging full URLs, and made jobs resumable with per-day
error handling. \\
\addlinespace

More data did not help &
SEC fundamentals and industry features lowered the development-period ranking
signal from IC~0.030 to between 0.007 and 0.027. &
They are not used in the 5-day model and are kept for the 1--12-month models
planned for spring. \\
\addlinespace

Kazakhstan data access &
The only free KASE price source is a broker API without published terms; the
official KASE API is paid. AIX has no free data, and some Kazakh news sites
block RSS. &
The prototype uses the broker API at a low request rate, for non-commercial use
only, and covers AIX companies through their KASE listings. Spring: request
academic access from KASE. \\
\addlinespace

Deployment and tooling &
nginx returned HTTP 502 after a container was recreated because it cached the
old IP address. The database port clashed with another service. CI had never
run because the repository was not on GitHub. &
nginx now resolves services on every request, and the database moved to port
55433. The repository was pushed on 26 November, and the first CI run passed in
3~min 19~s. \\

\end{longtable}
}

\section{Plans for Spring}
\label{sec:spring}

Spring adds the parts that need a GPU and fixes the evaluation problems found
this semester. The key result will be the price-only vs.\ price-plus-sentiment
ablation promised in the requirements document. Table~\ref{tab:spring} gives the
plan by month.

{\small
\begin{longtable}{@{}>{\raggedright\arraybackslash}p{0.11\linewidth}
                    >{\raggedright\arraybackslash}p{0.52\linewidth}
                    >{\raggedright\arraybackslash}p{0.30\linewidth}@{}}
\caption{Plan for the Spring 2027 semester.}
\label{tab:spring}\\
\toprule
\textbf{Month} & \textbf{Work} & \textbf{Done when} \\
\midrule
\endfirsthead
\multicolumn{3}{@{}l}{\small\emph{Table~\ref{tab:spring} continued}}\\
\toprule
\textbf{Month} & \textbf{Work} & \textbf{Done when} \\
\midrule
\endhead
\midrule
\multicolumn{3}{r@{}}{\small\emph{continued on the next page}}\\
\endfoot
\bottomrule
\endlastfoot

January &
Request NU HPC access and set up the RTX 5090 when it arrives. Build a
point-in-time stock universe from historical index membership. Download the
Alpaca news archive and add Finnhub or Alpha Vantage news. &
The GPU environment passes our verification script; the universe and the news
archive are in the data lake. \\
\addlinespace

February &
Score news sentiment in English, Russian and Kazakh. Compare FinBERT, a
zero-shot LLM API and a small open model on a hand-labelled test set. Run the
ablation price-only vs.\ price plus sentiment, on our data and on the StockNet
benchmark. &
An ablation table with accuracy, F1, MCC and Sharpe ratio for both data sets. \\
\addlinespace

March &
Fine-tune an open LLM with QLoRA on HPC. Add the stacking meta-learner and a
beta-neutral label. Connect the LLM trade reviewer to the sandbox, keeping the
rule-based reviewer as a fallback. &
The fine-tuned model beats the zero-shot baseline on held-out labelled news; the
reviewer returns valid verdicts for test trades in all three languages. \\
\addlinespace

April &
Build the Kazakhstan model, using next-session signals from London and NASDAQ
listings, Brent oil and USD/KZT. Train 1--12-month ranking models with the
fundamentals features already built. Request academic KASE data access, add a
nightly prediction job, and run user tests to improve the UI. &
KASE and long-horizon predictions appear on the website; user-test findings are
fixed. \\
\addlinespace

May &
Final evaluation on data not used for any design choice; final report and
presentation. &
Report and presentation submitted. \\

\end{longtable}
}
